% ============================================================================
% UPDATED ABSTRACT - CHILD LABOR PAPER
% January 2026
% ============================================================================

\begin{abstract}
\noindent This paper examines how supply chain monitoring affects child labor in developing countries, exploiting the 2013 Rana Plaza disaster as a quasi-experimental shock to factory oversight in Bangladesh. Using geocoded household survey data for 125,555 children linked to factory locations, we employ a difference-in-differences strategy comparing areas near textile factories to more distant locations. Standard estimates suggest a 1.7 percentage point reduction in child labor among children aged 10-13, the primary factory employment demographic. However, we document significant pre-treatment differential trends that complicate causal interpretation: treatment areas experienced 0.44 percentage points per year faster decline in child labor even before Rana Plaza. After controlling for group-specific trends, the treatment effect reverses sign, and using 2011 rather than 2007 as the reference year renders effects statistically insignificant. Event study analysis reveals that convergence between treatment and control areas began well before monitoring intensified. A formal sensitivity analysis following Rambachan and Roth (2023) confirms that results do not survive scrutiny: the breakdown value $\bar{M}^* = 0.00$ indicates estimated effects are insignificant even under exact parallel trends assumptions. We find no evidence of substitution to other work types or substantial increases in school enrollment among affected age groups. These findings highlight the importance of careful pre-trend analysis in difference-in-differences designs and suggest caution in attributing observed child labor declines to supply chain monitoring interventions. Our results contribute to debates about private governance effectiveness while underscoring methodological challenges in evaluating regulatory impacts absent randomized variation.

\vspace{0.5cm}
\noindent \textbf{JEL Classification:} J13, J81, O12, O14, F16, L67\\
\textbf{Keywords:} Child labor, Supply chain monitoring, Factory compliance, Rana Plaza, Bangladesh, Difference-in-differences, Pre-trends, Parallel trends assumption
\end{abstract}

% ============================================================================
% KEY NUMBERS FOR REFERENCE
% ============================================================================
%
% Sample: 125,555 children ages 5-17
% Treatment: 27,796 (22.1%) within 10km of factory
% Control: 97,759 (77.9%) beyond 10km
% Pre-period: 76,647 (61.0%)
% Post-period: 48,908 (39.0%)
% Overall child labor rate: 8.75%
%
% MAIN RESULTS (Ages 10-13):
% - Standard DiD: -1.66pp*** (SE = 0.0062)
% - With group-specific trend: +2.16pp** (SE = 0.0104)
% - Pre-treatment differential trend: -0.44pp/year*** (SE = 0.0009)
% - Reference year 2007: 2014 effect = -3.35pp***
% - Reference year 2011: 2014 effect = -0.81pp (n.s.)
%
% AGE HETEROGENEITY:
% - Ages 5-9: -0.16pp (n.s.), baseline CL = 0.6%
% - Ages 10-13: -1.64pp***, baseline CL = 7.1%
% - Ages 14-17: -0.07pp (n.s.), baseline CL = 21.1%
%
% ============================================================================
